
Consider the following experiment. Take a trained neural network and record its predicted test accuracy using a weight encoder. Now permute the neurons of a hidden layer --- shuffle which neuron is labeled ''channel 3'', ''channel 7'', and so on --- adjusting adjacent layers accordingly so the network computes exactly the same function. Re-encode and predict again. Repeat this across K = 50 such functionally equivalent rearrangements and average the predictions. For a correctly designed encoder, the average should match any individual prediction. For CISE, a sinusoidal positional encoder designed as a representative non-invariant baseline, the average yields $R^2 = -1.63$ --- substantially worse than predicting the mean ($R^2$ = 0). This result does not reflect robustness failure under noise. It reflects a precise measurement of a structural flaw: CISE encodes which channel occupies which position as a predictive signal. But channel positions in a neural network are arbitrary labels: two networks with identical function can differ only in which neuron happens to be labeled ''channel 3''. An encoder that treats these networks as different is not modeling what a network computes; it is modeling an arbitrary administrative choice made during initialization.

\subsubsection{1.1 The Weight-Space Symmetry Problem}

Neural networks with permutation-symmetric activation functions have a well-known symmetry: permuting the neurons of a hidden layer and correspondingly adjusting adjacent layers produces a functionally identical network \citep{HechtNielsen1990}. The space of weight configurations is therefore partitioned into orbits --- equivalence classes under permutation --- and any two configurations in the same orbit correspond to the same function.

Weight-space learning approaches, which train predictors directly on network weights, must contend with this symmetry. A predictor trained on one configuration may encounter a functionally identical configuration with different raw weights; if the encoder does not recognize these as equivalent, the predictor receives two different inputs that should produce the same output.

Prior work has addressed this from several angles. DeepSets \citep{Zaheer2017Deep} establishes that any permutation-invariant function on sets can be decomposed as $\rho$($\Sigma \phi (x_{i}))$, making sum pooling the canonical invariant architecture. Neural Functional Networks (NFN; Zhou et al., 2023) extend this to equivariant mappings over weight spaces. DWSNet \citep{Navon2023Equivariant} derives the complete set of affine equivariant and invariant linear layers for weight spaces and formalizes that the correct symmetry group for CNNs requires coupled row-column permutations across adjacent layers. Unterthiner et al. (2020) demonstrate that per-layer weight statistics achieve $R^2$ > 0.984 on CIFAR10-GS generalization prediction.

Despite this theoretical and empirical progress, a critical gap remains: no prior work has directly measured (1) how much non-invariant encoders vary in their representations of functionally equivalent networks (OrbitVar), (2) how much this representational variance propagates to prediction-space variance (MSE\_perm), and (3) whether architectural invariance causally closes this gap in downstream prediction quality.

\subsubsection{1.2 Our Approach}

We address this gap through a three-step experimental design that directly measures each link of the causal chain:

\textbf{Step 1 --- Architecture determines OrbitVar.} We measure within-orbit representational variance ($\text{OrbitVar} = \mathbb{E}_v[\text{Var}_{\pi}(\text{encoder}(\pi \cdot v))]$) for four encoders spanning the invariance spectrum: per-layer statistics (C0, approximately invariant), CISE sinusoidal positional encoder (C1, non-invariant), DeepSets doubly-invariant sum pooling (C2, exactly invariant), and NFN equivariant layers (C3, near-invariant).

\textbf{Step 2 --- OrbitVar propagates to prediction variance.} We apply a bias-variance decomposition over permutation orbits to measure how much of CISE's total prediction error is attributable to permutation-induced variance.

\textbf{Step 3 --- Eliminating MSE\_perm improves $R^2$.} We compare LightGBM predictors trained on each encoder's embeddings, with DeepSets' zero MSE\_perm as treatment and CISE's high MSE\_perm as control.

\subsubsection{1.3 Contributions}

This paper makes four contributions:

\textbf{1.} First joint measurement of OrbitVar, MSE\_perm, and $R^2$ across the full invariance spectrum on ModelZooDataset CIFAR10-GS. DeepSets achieves OrbitVar = 1.002e-14 (machine precision); NFN achieves 8.905e-08; CISE achieves 0.010333 --- a gap of 5 to 12 orders of magnitude, confirmed across 100 models with Wilcoxon p = 1.95e-18.

\textbf{2.} A novel MSE bias-variance decomposition over permutation orbits as a diagnostic tool for weight-space learning. For CISE, MSE\_perm/MSE\_total = 3.3452 --- permutation-induced variance contributes more than 3.35 times the total prediction error.

\textbf{3.} An empirical entanglement finding: the additive decomposition MSE = MSE\_res + MSE\_perm is empirically violated (closure = 0.872). MSE\_perm and MSE\_res are correlated in CISE embedding space, indicating that permutation sensitivity reshapes the downstream model's learned feature map rather than adding independent noise.

\textbf{4.} A practical encoder design result: DeepSets doubly-invariant encoding achieves $R^2$ = 0.9148 on ModelZooDataset CIFAR10-GS, a 6.37 percentage-point improvement over CISE ($R^2$ = 0.8511), with no additional training supervision.

